\chapter{Introduction}
\label{introduction}

Humans navigate complex environments by leveraging their perceptual capabilities in order to act purposefully in existential terms. Upon making an observation they must consider arbitrarily many causes in the environment. Knowing the causes of observations is important because they may pose a threat which the indiviual should counteract via actions. Having to account for a possibly indefinite number of origins of observations means a computational challenge for the human brain. Evolution has come up with a solution to this: Abstraction. Instead of having a model for each and every detail in the environment, humans are able to abstract and thus can file away functionally similar causes as singular entities. For example, when a person has learned how to do grocery shopping at a certain store, then she is also able to this at another store without relearning every detail of the new shop and how to navigate it. She posseses the abstract knowledge of grocery shopping and can apply this knowledge purposefully even if the details have changed.

Abstraction can be considered as having a context state that is "activated" during certain situations. In this way, perception is interpreted in the light of a specific context. This has been investigated in different forms: \textcite{millerandcohen2001} proposed that the prefrontal cortex maintains certain (context) patterns of activity that provide bias signals to other brain structures. In this case these patterns represent goals of the individual. With respect to vision, \textcite{raoandballard1999} discussed a hierarchy of cortical areas where higher-order areas provide (contextual) feedback connections to lower-order areas to efficiently process visual input. Moreover, the balancing between habitual- and goal-directed behavior was also proposed with respect to a context state: \textcite{schwoebel2021} developed a computational model of behavior in which an abstract context state (to be inferred) then either leads to automatic habitual behavior or starts a computationally demanding planning process in which different sequences of actions are evaluated. This model was discussed with respect to substance use disorder, which is a habitual trait that is context dependent. Concluding, the existence of abstract context states is evident both in neural wetware and in models of behavior.

As touched upon earlier, existence is the central problem every human must face. While contextual abstraction allows for tractable computation (inference) of causes (that generate observations) the mechanisms of utilizing this knowledge of hidden environmental causes remains to be addressed: According to the free energy principle (FEP) (\cite{friston2010}), individuals are equipped with a generative model that relates environmental causes with observations. Moreover, it recognizes that existence can be equated with maintaining a non-equilibrium steady state (NESS) distribution (\cite{friston2013}), i.e, the individual "wants" to remain in a state that is typical for its ecological niche and phenotype. In other words, the animal is inclined to maintain likely states that underwrites its existence. In its more general form the principle suggests that the whole individual is a generative model. Under this perspective the individual might be seen as a functional homomorphism of its ecological niche, i.e, it keeps stable (survives) because environmental influences (including observations) get "explained away" by its physiological structure. In the FEP it is then concluded that animals maintain their NESS by minimizing the variational free energy (VFE) with respect to their generative models. Moreover, VFE is a quantity that relates to the present because existence itself relates to the present. Following this, the only self consistent belief an individual can have is that it also minimizes its VFE in the future, i.e., it should perform actions that fulfill this requirement. Interestingly, when planning sequences of actions that minimize VFE in future the animal does not solely realize preferred observations that are typical for its ecological niche (e.g., most humans do not "like" the gaze of spiders because they are potentially venomous) but also shows curiousity (\cite{friston2015}), i.e., it seeks information in the environment that disambiguates hidden causes (of observations) which makes its life more predictable and therefore secure.

Following this, it is plausible that an individual not only shows curiousity with respect to immediate environmental causes but also attempts to seek information about the more abstract context that might currently be active. For example, consider a group of people who plays cards. A new player likely asks for the rules before joining. In this situation the new player wants to find out about the context "game rules" by seeking information in the environment "asking the other players". This pattern seems to be cental to human behavior.

In this thesis I attempt to develop a behavioral computational model that may account for contextual information seeking in humans. The model integrates central pillars of intelligent behavior, e.g., perception, planning, action, and inference. A central motivation for this work is the model's usage in Bayesian model comparison, i.e, it can be understood as a hypothesis about behavior that can be validated by empirical data from experiments with human participants. While this perspective makes it useful for neuroscientific and psychological research, in this work I view the behavioral model as an agent to be able to explicitly show the kind of behavior that can be explained by this model. Concretely, I develop an agent and demonstrate its contextual information seeking behavior in simulated experiments in the T-Maze envirnonment that is usually used to show this kind of behavior in discrete agents that function akin to the free energy principle (\cite{friston2015}). This last point effectively makes the agent an active inference agent (\cite{friston2010aif}).

Active inference agents are behaving entities that entail generative models that allow them to "understand" the relationships between observations, hidden environmental states (causes), and the succession of hidden states when influenced by actions. Using this model, the agent can integrate current observations, infer its likely causes in the environment, simulate how sequences of actions lead to future observations that it either likes or dislikes (with respect to its preferences), and finally select actions that it expects to yield the lowest variational free energy in the future. These behavioral "functions" can be framed as a unified inference process during which the agent infers belief distributions with respect to observations, hidden states, and actions at different timesteps (past, present, and future). When developing computational models of these mechanisms, Forney factor graphs have proven to be useful theoretical substrates that are able to capture the generative model (unrolled in time), the belief distributions, and the unified inference process as messages sent through this graph (\cite{loeliger2004,delaar2019}). While the usage of Forney graphs makes the underlying mathematical foundation of active inference agents more visually accessible they also lend some plausability as a model of brain function where computations are equivalently distributed and neural signaling might be analog to local messages sent in the factor graph (\cite{bastos2012,friston2017}). Following this, the agent to be developed in this work will be based on a Forney factor graph representation.

Specifically, I will resort to an approach presented by \textcite{senoez2021,delaar2023,delaar2025}. In this, we start with the definition of the agent's generative model (unrolled in time). Next, we view this model as its Forney factor graph representation and introduce a Bethe free energy \(F\) that relates different parts of the generative model with additional belief distributions \(q\) that exist on the edges (and nodes) in the factor graph. Following this, we introduce constraints to this free energy that relate to local marginalization constraints, but also to the information seeking mechanism among others, which leads to a corresponding Lagrangian \(L\). From this we derive stationary solutions for the belief distributions and leverage the marginalization constraints to derive local messages in the factor graph. During this derivation it also becomes obvious that these messages allow for a direct computation of the belief distributions. Effectively, I derive a message passing algorithm for the agent that implements an inference process and yields infered belief distributions as a result.

Towards the development of an active inference agent that supports contextual information seeking (named Context-GFE agent in this work), I also implement three other agents that serve as starting points for this attempt: First, I reproduce two simpler agents (BFE- and GFE agent) from ... that do not entail an abstract context state but only (normal) hidden states as part of their generative models. Both are similar but the GFE agent additionally is able to perform information seeking over its (normal) hidden state. In this work, the introduction of these two agents has three reasons: First, I can show how the integration of the information seeking mechanism in the GFE agent over the BFE agent generally works. Second, both agents can serve as baselines during the simulations to reveal the kind of behavior that is expected from an agent that performs informations seeking (GFE agent) in comparison to an agent without this ability (BFE agent). Third, the BFE agent serves as a basis towards the implementation of the Context-GFE agent. This last point requires further explanation: I expand the BFE agent by an additional context state which leads to the Context-BFE agent. As a side note: The Context-BFE agent is also inspired by the Bayesian contextual control model developed by \textcite{schwoebel2021} but lacks the balancing between habitual and goal-directed behavior. Having developed the Context-BFE agent I finally take the integration of the information seeking mechanism in the GFE agent as a blueprint to implement this mechanism in the Context-BFE agent–-–but over its context state, which eventually leads to the Context-GFE agent.

In \cref{sec:0} I give an introduction to active inference. \cref{sec:1} presents the general procedere of deriving message passing algorithms from custom Lagrangians that relate generative models with belief distributions represented by Forney factor graphs. Moreover, in \cref{sec:2} I derive the four agents (BFE, GFE, Context-BFE, and Context-GFE) and present the experimental setups in \cref{sec:3}. \cref{results} presents the results of the experiments and in \cref{discussion} I give a corresponding discussion. The thesis concludes with \cref{conclusion}.